\documentclass[a4paper, 12pt, hyphens]{book}

% Pacchetti utili
% Codifica
\usepackage[utf8]{inputenc}

% Lingua, sostituire italian con english nel caso in cui la tesi sia scritta in inglese
\usepackage[italian]{babel}
\usepackage[T1]{fontenc}

% Stile font
\usepackage{lmodern}

% Colori
\usepackage[usenames]{color}
\usepackage[dvipsnames]{xcolor}
\usepackage{colortbl}

% Pacchetto per definire layout di pagina
%\usepackage{sectsty}
\usepackage{fancyhdr}
\usepackage[left=3cm, right=3cm, bottom=3cm]{geometry}

% Spazia linee all'interno del documento
\usepackage{setspace}

% Impostazioni note a piè pagina
\usepackage[stable, multiple]{footmisc}

% Crea link ipertestuali
\usepackage[hidelinks, hyperfootnotes=false]{hyperref}

% PDF in formato PDF/A-1b
%\usepackage[a-1b]{pdfx}

% Inclusione immagini
\usepackage{graphicx}

% Didascalie immagini
\usepackage[hang, small, sf, labelfont=bf]{caption}
\usepackage{subcaption}

% Formattazione url
\usepackage{url}

% Inserimento formule
\usepackage{amsmath}
\usepackage{mathrsfs}

% Listati di codice
\usepackage{verbatim}
\usepackage{listings}

%Inserimento pseudocodice
\usepackage{algorithm}
\usepackage{algpseudocode}

% Citazione frasi
\usepackage{csquotes}

% Citazioni e riferimenti a label
%\usepackage{cite} % Don't use with biblatex package
%\usepackage[english]{varioref}

% Bibliografia
% This ----------------------------------------------------------
\usepackage[backend=biber, style=numeric, sorting=none, backref=true]{biblatex} % biber or bibtex - numeric or alphabetic - none or nty
\addbibresource{chapters/references.bib}
% Or this -------------------------------------------------------
% \usepackage{natbib} 
% \usepackage[nottoc]{tocbibind}
% \bibliographystyle{plain}

\usepackage{chngcntr}
% \counterwithin{figure}{section} % Change figure numbering by section (with chngcntr pkg)
% \counterwithin{table}{section} % Change table numbering by section (with chngcntr pkg)
% \numberwithin{equation}{section} % Change equation numbering by section (with amsmath pkg)
\counterwithout{footnote}{chapter} % Change footnote numbering by chapter (with chngcntr pkg)

% Inserimento di Lorem ipsum nel testo
\usepackage{lipsum}

% Prima pagina bianca
\usepackage{afterpage}
\newcommand\myemptypage{
                        \null
                        \thispagestyle{empty}
                        \addtocounter{page}{-1}
                        \newpage
                        }

\usepackage{tabularx} 

%\usepackage{marginnote}

% Impostazioni dei margini, definizione di colori e stili
\newfont{\ttvar}{cmvtt10 scaled 1200}   % nuovo carattere tipi courier a spaziatura variabile per le dimostrazioni

% margini senato
\textwidth       =  14.50 cm            % larghezza 21 cm - 4 cm (sinistro) - 2.5 (destro)
\textheight      =  23.10 cm            % altezza 29.7 cm - 3 cm (superiore) - 2 (inferiore)
\topmargin       =   0.00 cm            % margine superiore 3 cm diminuito di 1 inch
\oddsidemargin   =   1.46 cm            % margine sinistro 4 cm diminuito di 1 inch
\evensidemargin  =  -0.04 cm            % margine destro 2.5 cm diminuito di un inch

\setlength{\headsep}{1.0cm}
\setlength{\footskip}{1.0cm}
\setlength{\headheight}{15pt} % ...at least 14.9999pt
\parindent = 0.7cm
\captionmargin = 0.7cm

%Interlinea 1.5
\onehalfspacing

% stile pagina
\pagestyle{fancy}
\renewcommand{\chaptermark}[1]{\markboth{\chaptername\ \thechapter.\ #1 }{}}
\renewcommand{\sectionmark}[1]{\markright{\thesection\ #1}{}}
\fancyhead{}
\fancyhead[LE,RO]{\sffamily \thepage}
\fancyhead[RE]{\sffamily \leftmark}
\fancyhead[LO]{\sffamily \rightmark}
\fancyfoot{}%[CE,CO]{\sffamily \thepage}

%\renewcommand{\headrulewidth}{0pt}
% ridefinisco lo stile plain
\fancypagestyle{plain}{ \fancyhead{} \fancyfoot{}
                        \fancyfoot[C]{\sffamily \thepage}
                        \renewcommand{\headrulewidth}{0pt}
                      }

% stile per i titoli
%\allsectionsfont{\sffamily \raggedright}

%definisco stile listati di codice
\lstdefinestyle{mystyle}{
                        backgroundcolor=\color[RGB]{240,240,240}, %Red, Green, Blue from 0 to 255
                        basicstyle=\color{black}\ttfamily\footnotesize,
                        commentstyle=\color{OliveGreen}\textit\tiny,
                        keywordstyle=\color{blue},
                        identifierstyle=\color{black},
                        stringstyle=\color{violet},
                        breakatwhitespace=false,         
                        breaklines=true,
                        captionpos=b, %bottom, top
                        rulecolor=\color{black},
                        frame=trBL, %Top,Right,Bottom,Left (Case sensitive)
                        frameround=tttt, %t=round corner, f=straight corner
                        keepspaces=true,                 
                        numbers=left,                    
                        numbersep=10pt,
                        numberstyle=\color{gray},
                        showspaces=false,                
                        showstringspaces=false,
                        showtabs=false,                  
                        tabsize=1,
                        columns=fullflexible %makes code copy and pastable
                        }

\lstset{style=mystyle}
\lstset{emph={RandomForestClassifier, RandomizedSearchCV, GridSearchCV},emphstyle=\underbar}

\renewcommand{\lstlistingname}{Codice} % Per rinominare la descrizione dei codici
\renewcommand{\lstlistlistingname}{Elenco dei codici} % Per rinominare la descrizione dei codici

% Nuovi comandi
% \newcommand{\abs}[1]{\left|#1\right|}                               % modulo
\newcommand{\dato}{\left|\right.}                                   % probabilit\`{a} condizionata
\newcommand{\fun}[1]{\mathrm{#1}}                                   % stile funzione
\newcommand{\imp}{\;\;\Longrightarrow\;\;}                          % implicazione
\newcommand{\norma}[1]{\left\| #1 \right\|}                         % norma
\newcommand{\prob}[1]{\mathrm{P}\!\left[#1\right]}                  % probabilit\`{a}
\newcommand{\expect}[1]{\mathrm{E}\!\left[#1\right]}                % aspettazione
\newcommand{\sse}{\;\;\Longleftrightarrow\;\;}                      % se e solo se
\newcommand{\vect}[1]{{\boldsymbol{\mathrm{#1}}}}                   % stile vettore
\newcommand{\real}[1]{\fun{Re}\left[#1\right]}                      % parte reale
\newcommand{\imag}[1]{\fun{Im}\left[#1\right]}                      % parte immaginaria
\newcommand{\Dim}[1]{\fun{dim}\left[#1\right]}                      % dimensione di una matrice
\newcommand{\Det}[1]{\fun{det}\left[#1\right]}                      % determinante di una matrice
\newcommand{\Ker}[1]{\fun{ker}\left[#1\right]}                      % ker di una matrice
\newcommand{\rango}[1]{\fun{rango}\left[#1\right]}                  % rango di una matrice
\newcommand{\scalare}[2]{\left\langle #1, #2 \right\rangle}         % prodotto scalare
\newcommand{\blbrace}{\left  \lbrace}                               % parentesi graffa sinistra grande
\newcommand{\brbrace}{\right \rbrace}                               % parentesi graffa destra grande
\newcommand{\sinc}{\fun{sinc}}                                      % sinc
\newcommand{\rect}{\fun{rect}}                                      % rect
\newcommand{\rcos}{\fun{rcos}}                                      % rcos
\newcommand{\sgn}{\fun{sgn}}                                        % sgn
\newcommand{\N}{\mathbb{N}}                                         % insieme dei numeri naturali
\newcommand{\Z}{\mathbb{Z}}                                         % insieme dei numeri interi
\newcommand{\Q}{\mathbb{Q}}                                         % insieme dei numeri razionali
\newcommand{\R}{\mathbb{R}}                                         % insieme dei numeri reali
%\newcommand{\C}{\mathbb{C}}                                         % insieme dei numeri complessi
\newcommand{\seq}[2][n]{#2_{0}, #2_{1}, \ldots, \, #2_{#1}}         % sequenza
\newcommand{\Span}[2][n]
{\fun{span} \blbrace #2_{1}, #2_{2}, \ldots, \, #2_{#1} \brbrace}   % spazio generato
\newcommand{\ddt}{\frac{\fun{d}}{\fun{dt}}}                         % derivata
\newcommand{\Div}[2]{#1 \; \mid \; #2}                              % divide
\newcommand{\MCD}[2]{\fun{MCD}\(#1, #2\)}                           % massimo comun divisore
\newcommand{\mcm}[2]{\fun{mcm}\(#1, #2\)}                           % minimo comune multiplo
\newcommand{\goodgap}{
                      \hspace{\subfigtopskip}
                      \hspace{\subfigbottomskip}
                     }                                              % interlinea opportuna per le sottofigure
\newcommand{\eng}[1]{\emph{#1}}                                   % inglese
\newcommand{\virg}[1]{``#1"}                                        % fa una citazione tra virgolette
%\newcommand{\unit}[2]($\frac{\text{#1}}{\text{#2}}$)                % unit\`{a} di misura
\newcommand{\textttvar}[1]{{\ttvar #1}}

\definecolor{gray}{gray}{0.9}
\newcommand{\listato}[1]{\lstset{
                                language=#1,
                                numbers=left,
                                numberstyle=\tiny, 
                                stepnumber=2, 
                                numbersep=5pt, 
                                numberblanklines=false, 
                                xleftmargin=5pt, 
                                captionpos=b, 
                                stringstyle=\ttfamily, 
                                columns=flexible, 
                                showstringspaces=false, 
                                tabsize=2, 
                                frame=single, 
                                framerule=0pt, 
                                backgroundcolor=\color{gray}, 
                                basicstyle=\small}
                                }

%****************************
%* ridefinizioni di comandi *
%****************************

\renewcommand{\(}{\left(}                                     % parentesi tonda sinistra grande
\renewcommand{\)}{\right)}                                    % parentesi tonda destra grande
\renewcommand{\[}{\left[}                                     % parentesi quadra sinistra grande
\renewcommand{\]}{\right]}                                    % parentesi quadra destra grande
\renewcommand{\exp}[1]{\fun{e}^{#1}}                          % esponenziale
\renewcommand{\gcd}[2]{\fun{gcd}\(#1, #2\)}                   % massimo comun divisore

\renewcommand{\lstlistingname}{Listing}
\renewcommand{\lstlistlistingname}{Elenco dei listati codice}

\makeatletter
\def\BState{\State\hskip-\ALG@thistlm}
\makeatother

%Definizione input e output pseudocodice
\renewcommand{\algorithmicrequire}{\textbf{Input:}}
\renewcommand{\algorithmicensure}{\textbf{Output:}}
 
%"Nome" della bibliografia
\addto{\captionsenglish}{%
  \renewcommand{\bibname}{References}
}

\begin{document}

    \frontmatter

    \myemptypage % Pagina bianca prima del titolo
    
    % Prima pagina ----------------------------------------------------------------------

    \begin{titlepage}

    \begin{center}
        \includegraphics[scale=0.1]{images/titlepage/logo_unipd_small}
        
        \vspace{1cm}
        
        \textsc{\LARGE Universit\`{a} degli Studi di Padova}
        
        \vspace{0.5cm}
        
        \textsc{\large Dipartimento di ...}
        
        \vspace{0.5cm}
        
        \textsc{\large Corso di Laurea ...\\in Ingegneria ...}
        
        \vfill
        
        {\Huge \bfseries Titolo tesi}
        
        \vspace{1cm}
        
        {\LARGE Titolo inglese}
    \end{center}
    
    \vfill
    
    \noindent \textit{\large Relatore:} \\
    \textsc{\large Ch.mo Prof. Nome Cognome}
    %\textsc{\large Chiar.ma Prof.ssa Nome Cognome}
    
    %\vfill
    \vspace{1cm}
    
    \begin{flushright}
        \textit{\large Laureando:}   \\
        \textsc{\large Nome Cognome} \\
        \textsc{1234567}
    \end{flushright}
    
    \vfill
    
    \centering {\large Anno Accademico 20XX/20XX}
    
\end{titlepage}
    
    \thispagestyle{empty} % Pagina bianca dopo il titolo
    \cleardoublepage
    
    % Citazione, ringraziamenti, abstract, indice ---------------------------------------
    
    \mainmatter  
    
    \pagenumbering{roman} % Numerazione romana
    \thispagestyle{empty}
    
    \clearpage{\pagestyle{plain}\cleardoublepage} % Comando per iniziare il capitolo su pagina dispari
    % Citazione
\vspace*{\fill} 
\begin{quote} 
    \centering 
    \begin{flushright}{
                        \slshape  
                        Inserisci qui la tua citazione o dedica.
                      }
    \end{flushright}
\end{quote}
\vspace*{\fill}

\clearpage{\pagestyle{plain}\cleardoublepage}

% Ringraziamenti
\chapter*{Ringraziamenti}
Inserisci gli eventuali ringraziamenti personali.
	
    
    \clearpage{\pagestyle{plain}\cleardoublepage} 
    % %definisco il layout dell'abstract
% \def\changemargin#1#2{\list{}{\rightmargin#2\leftmargin#1}\item[]}
% \let\endchangemargin=\endlist 

% %Genero l'ambiente per l'abstract
% \newcommand\summaryname{Sommario}
% \newenvironment{Abstract}%
%     {\begin{center}%
%     \bfseries{\summaryname} \end{center}}
    
% \begin{Abstract}
%     \begin{changemargin}{1cm}{1cm}
%     Inserire abstract. I margini nell'abstract sono stati ridotti di un centimetro. In caso non si volesse questa riduzione rimuovere \textit{changemargin}.
%     \end{changemargin}
% \end{Abstract}

\chapter*{Sommario}
\label{sommario}

Inserire abstract.
    
    \clearpage{\pagestyle{plain}\cleardoublepage} 
    \chapter*{Introduzione}
\label{introduzione}

Scrivi qui la tua introduzione.
    
    \clearpage{\pagestyle{plain}\cleardoublepage}
    \tableofcontents % Indice
    
    \clearpage{\pagestyle{plain}\cleardoublepage}
    \phantomsection
    \addcontentsline{toc}{chapter}{Elenco delle figure} % Per inserire nell'indice la lista dei codici (non è automatico)
    \listoffigures % Lista figure
    
    \clearpage{\pagestyle{plain}\cleardoublepage}
    \phantomsection
    \addcontentsline{toc}{chapter}{Elenco delle tabelle} % Per inserire nell'indice la lista dei codici (non è automatico)
    \listoftables % Lista tabelle
    
    \clearpage{\pagestyle{plain}\cleardoublepage}
    \phantomsection
    \addcontentsline{toc}{chapter}{Elenco dei codici} % Per inserire nell'indice la lista dei codici (non è automatico)
    \lstlistoflistings % Lista codici
    
    % Capitoli --------------------------------------------------------------------------
    
    \clearpage{\pagestyle{plain}\cleardoublepage}
    \pagenumbering{arabic} % Numerazione araba per i capitoli
    
    \clearpage{\pagestyle{plain}\cleardoublepage} 
    \chapter{Primo Capitolo} % Nome capitolo
\label{chapter:primo_capitolo} % Label per creare riferimenti al capitolo
La struttura utilizzata in questo template non è obbligatoria, però ritengo che sia molto comoda per evitare di scrivere file troppo lunghi e di avere un controllo migliore sulla struttura. Questa prevede di scrivere l'introduzione al capitolo in un file salvato nella cartella principale e di sviluppare le sezioni all'interno di una cartella. Esempio di richiamo ad un riferimento \cite{Vas}.

\section{Sezione 1}
Ad ogni sezione, in questo template \cite{UsabilityFirst}, corrisponde un file all'interno della cartella relativa al capitolo \cite{Bolognani-Zigliotto}.

\section{Sezione 2}
Una sezione può contenere una sottosezione. In questo caso, è stato deciso di non creare altri file... \cite{mori:poesie}

\subsection{Sottosezione}
...ma di scrivere il testo all'interno del file della sezione corrente \cite{Munoz-Lipo}.  % File in cui verrà scritto il capitolo
    
    \clearpage{\pagestyle{plain}\cleardoublepage}
    \chapter{Immagini e Tabelle} 
\label{chapter:immagini_e_tabelle} 
Metodi più comuni per inserire immagini e esempi di tabelle. \cite{eco:tesi}

\section{Immagine singola}

Per fare riferimento ad un immagine \cite{dirac}, come ad ogni altro elemento a cui viene attribuita una \textit{label} è disponibile il comando \textit{ref}. Riferimento a Figura \ref{fig:figura_test}.
\begin{figure}[!ht]
    \centering
    \includegraphics[width=0.9\textwidth]{images/cap2/test}
    \caption{Didascalia}
    \label{fig:figura_test}
\end{figure}

\section{Immagine multipla}

Inserire più ``sottofigure" in una figura.
\begin{figure}[!ht]
    \centering
    
	%Prima riga
	\begin{subfigure}{0.49\textwidth} %Controllo della posizione in orizzontale
		\includegraphics[width=0.9\textwidth]{images/cap2/test} 
		\caption{}
	\end{subfigure}
	\begin{subfigure}{0.49\textwidth}
		\includegraphics[width=0.9\textwidth]{images/cap2/test}
		\caption{}
	\end{subfigure}
	
	%Seconda riga
	\begin{subfigure}{0.49\textwidth}
		\includegraphics[width=0.9\textwidth]{images/cap2/test} 
		\caption{}
	\end{subfigure}
	\begin{subfigure}{0.49\textwidth}
		\includegraphics[width=0.9\textwidth]{images/cap2/test}
		\caption{}
	\end{subfigure}
    \caption{Esempio di figura composta da 4 figure.}
\end{figure}

%\cleardoublepage

\section{Tabelle}

Nella seguente tabella vengono mostrati alcuni esempi di separazione delle righe. La separazione delle colonne avviene  all'interno delle parentesi graffe dopo il comando \textit{tabular}.
\begin{table}[h]
    \centering
    \begin{tabular}{ ||c|c|c|| } 
        \hline
        cella1  & cella2  & cella3  \\
        \hline
        \hline 
        cella4  & cella5  & cella6  \\ 
        \hline
        cella7  & cella8  & cella9  \\ 
        cella10 & cella11 & cella12 \\ 
        \hline
    \end{tabular}
    \caption{Tabella 1}
    \label{tab:tabella1}
\end{table}

È possibile, inoltre, fissare la dimensione delle colonne \cite{prova2}.
\begin{table}[ht]
    \centering
    \begin{tabular}{ | m{3cm} | m{5cm} | m{1.5cm} | } 
        \hline
        cella1  & cella2  & cella3                      \\
        \hline
        \hline 
        cella4  & cella5  & cella6 cella6 cella6 cella6 \\ 
        \hline
        cella7  & cella8  & cella9                      \\ 
        cella10 & cella11 & cella12                     \\ 
        \hline
    \end{tabular}
    \caption{Tabella 2}
    \label{tab:tabella2}
\end{table} 
    
    \clearpage{\pagestyle{plain}\cleardoublepage}
    \chapter{Formule} 
\label{chapter:formule}
Per inserire delle formule matematiche è possibile utilizzare due metodi \cite{prova1}:
\begin{itemize}
    \item in linea: inserendo la formula tra due caratteri \$. 
    \item utilizzando l'ambiente \textit{equation}.
\end{itemize}

%Il comando noindent viene utilizzato per evitare l'indentazione del paragrafo
\noindent Esempio di formula in linea $x(t)=x_{0} + v_{0}t + \frac{1}{2}at^{2}$.

\noindent Esempio di utilizzo dell'ambiente \textit{equation}:
\begin{equation}\label{eq:legge_oraria}
    x(t)=x_{0} + v_{0}t + \frac{1}{2}at^{2}
\end{equation}
Possiamo usare le \textit{label} anche per le equazioni \cite{wiki:latex}. Legge oraria nell'Equazione \ref{eq:legge_oraria}.
Infine, un esempio di formula su più righe:
\begin{equation}
    \begin{split}
    x(t) & = x_{0} + v_{0}t + \frac{1}{2}at^{2} \\
    & =	x_{0} + v_{0}t + \frac{1}{2}\frac{F}{m}t^{2}
    \end{split}
\end{equation} 
    
    \clearpage{\pagestyle{plain}\cleardoublepage}
    \chapter{Pseudocodice e codice} 
\label{chapter:codice} 
In questo template per l'inserimento di pseudocodice è stato utilizzato il pacchetto \textit{algpseudocode}. Per quanto riguarda l'inserimento del codice è possibile utilizzare il comando \textit{verb} per inserire in linea oppure \textit{lstlisting} per inserire blocchi di codice \cite{einstein}. 

\section{Pseudocodice}

\begin{algorithm}
    \caption{Nome algortimo}\label{algo:name}
    \begin{algorithmic}[1]
        \Require $\text{Input dell'algoritmo}$
        \Ensure $\text{Output dell'algortimo}$
        \State variabile $\gets$ assegnazione valore
        \State valore\_ritornato $\gets$ \textsc{Funzione}\textsc{Prova}(param1, param2)
        
        \For{\textit{element} in \textit{list}} 
        \State res $\gets$ \textsc{Do}\textsc{Something}(element)
        \EndFor
        
        \If{condizione1}
        \State do something
        
        \ElsIf{condizione2}
        \State do something else
        \Else
        \State \textbf{print} ``Hello World"
        \EndIf
    \end{algorithmic}
\end{algorithm}

\section{Codice}

È possibile inserire codice in linea: \verb!print("Hello World")! \cite{knuthwebsite}.\\
Inoltre è possibile usare l'ambiente \textit{lstlisting} configurando il layout del blocco di codice nel file \textit{layout.tex}. Il codice può essere importato da un file esterno che metteremo nella cartella \textit{code} \cite{mori:tesi}.
\lstinputlisting[   language=Python,
                    caption=Didascalia.,
                    label=ls:codice]
                    {code/randomized_search.py}
            
\lstinputlisting[
                    language=C++,
                    %linerange={0-22},
                    %firstnumber=20,
                    %frame=trL,
                    % frameround=tttt,
                    caption=prova.cpp,
                    label=list:prova
                    ]
                    {code/prova.cpp}
    
    \clearpage{\pagestyle{plain}\cleardoublepage}
    \chapter*{Conclusioni} \addcontentsline{toc}{chapter}{Conclusioni}
\label{conclusioni}

Scrivi qui le tue conclusioni.
    
    % Bibliografia ----------------------------------------------------------------------
    
    \clearpage{\pagestyle{plain}\cleardoublepage}
    
    % This -------------------------------
    %\printbibliography[heading=bibintoc]
    
    % Or this ----------------------------
    % \chapter*{Bibliografia} \addcontentsline{toc}{chapter}{Bibliografia}
    % \printbibliography[heading=subbibintoc, type=book,    title={Libri}]
    % \printbibliography[heading=subbibintoc, type=article, title={Articoli}]
    % \printbibliography[heading=subbibintoc, type=online,  title={Siti}]
    % \printbibliography[heading=subbibintoc, type=misc,    title={Vari}]
    
    % Or this ----------------------------
    \chapter*{Bibliografia} \addcontentsline{toc}{chapter}{Bibliografia} \fancyhead[RE,LO]{\sffamily Bibliografia}
    \printbibliography[nottype=online, heading=subbibliography, keyword={cap1}, title={Capitolo 1}]
    \printbibliography[nottype=online, heading=subbibliography, keyword={cap2}, title={Capitolo 2}]
    \printbibliography[nottype=online, heading=subbibliography, keyword={cap3}, title={Capitolo 3}]
    \printbibliography[nottype=online, heading=subbibliography, keyword={cap4}, title={Capitolo 4}]
    
    \chapter*{Sitografia} \addcontentsline{toc}{chapter}{Sitografia} \fancyhead[RE,LO]{\sffamily Sitografia}
    \printbibliography[type=online, heading=subbibliography, keyword={cap1}, title={Capitolo 1}]
    \printbibliography[type=online, heading=subbibliography, keyword={cap2}, title={Capitolo 2}]
    \printbibliography[type=online, heading=subbibliography, keyword={cap3}, title={Capitolo 3}]
    \printbibliography[type=online, heading=subbibliography, keyword={cap4}, title={Capitolo 4}]
        
\end{document}
